\documentclass[10pt,a4paper]{article}
\usepackage[margin=20mm,headheight=15pt,headsep=7mm,footskip=10mm]{geometry}
\usepackage{fontspec,amsmath,amssymb,mathtools,microtype}
\setmainfont{lmroman10-regular.otf}[BoldFont=lmroman10-bold.otf,ItalicFont=lmroman10-italic.otf,BoldItalicFont=lmroman10-bolditalic.otf]
\setsansfont{lmsans10-regular.otf}[BoldFont=lmsans10-bold.otf,ItalicFont=lmsans10-oblique.otf,BoldItalicFont=lmsans10-boldoblique.otf]
\setmonofont[Scale=.86]{lmmono10-regular.otf}[ItalicFont=lmmono10-italic.otf]
\usepackage{ragged2e}\usepackage{xcolor,booktabs,longtable,array,graphicx,enumitem,fancyhdr,needspace,titlesec,hyperref,xurl}
\definecolor{AtlasNavy}{HTML}{18334C}\definecolor{AtlasTeal}{HTML}{116D73}\definecolor{AtlasGray}{HTML}{52616D}
\hypersetup{colorlinks=true,linkcolor=AtlasTeal,urlcolor=AtlasTeal,pdfcreator={Complexity Atlas},unicode=true}
\urlstyle{same}\setlength{\parindent}{0pt}\setlength{\parskip}{4.8pt}\setlength{\emergencystretch}{3em}
\renewcommand{\baselinestretch}{1.08}\setcounter{secnumdepth}{0}\setlist{nosep,leftmargin=1.5em}
\titleformat{\section}{\Large\bfseries\sffamily\color{AtlasNavy}\raggedright}{}{0em}{}
\titleformat{\subsection}{\large\bfseries\sffamily\color{AtlasTeal}\raggedright}{}{0em}{}
\titleformat{\subsubsection}{\normalsize\bfseries\color{AtlasNavy}\raggedright}{}{0em}{}
\titlespacing*{\section}{0pt}{14pt}{6pt}\titlespacing*{\subsection}{0pt}{10pt}{5pt}
\pagestyle{fancy}\fancyhf{}\fancyhead[L]{\small\sffamily\color{AtlasGray}COMPLEXITY ATLAS}
\fancyhead[R]{\small\sffamily\color{AtlasGray}ENGLISH EDITION}\fancyfoot[C]{\small\thepage}\renewcommand{\headrulewidth}{.3pt}
\newcommand{\AtlasMath}[1]{\par\begingroup\centering\setbox0=\hbox{$\displaystyle #1$}\ifdim\wd0>1.4\linewidth\typeout{ATLAS WIDE FORMULA: \the\wd0}\fi\ifdim\wd0>\linewidth\resizebox{\linewidth}{!}{\box0}\else\box0\fi\par\endgroup}
\newcommand{\AtlasInlineMath}[1]{\begingroup\setbox0=\hbox{$#1$}\ifdim\wd0>.4\linewidth\endgroup\AtlasMath{#1}\else\endgroup$#1$\fi}
\pdfstringdefDisableCommands{\def\AtlasInlineMath#1{#1}}
\DeclareMathOperator{\polylog}{polylog}\DeclareMathOperator{\dist}{dist}\DeclareMathOperator{\diam}{diam}
\DeclareMathOperator{\rank}{rank}\DeclareMathOperator{\tr}{tr}\DeclareMathOperator{\diag}{diag}
\DeclareMathOperator{\prox}{prox}\DeclareMathOperator{\KL}{KL}\DeclareMathOperator{\Var}{Var}
\providecommand{\eps}{\varepsilon}\providecommand{\R}{\mathbb R}\providecommand{\E}{\mathbb E}
\providecommand{\mathbbm}[1]{\mathbf{#1}}\providecommand{\mathds}[1]{\mathbf{#1}}
\providecommand{\llbracket}{[\![}\providecommand{\rrbracket}{]\!]}
\begin{document}
\raggedright

{\LARGE\bfseries\sffamily\color{AtlasNavy} Inequality-constrained lower level: primitive information and Goldstein stationarity\par}\vspace{5pt}
\textbf{Setting ID.} \texttt{BL\_\allowbreak{}SPEC\_\allowbreak{}linear\_\allowbreak{}inequality}

Bilevel optimization

Let \AtlasInlineMath{h(x,y)=Ax-By-b}, \AtlasInlineMath{y^*(x)=\operatorname*{argmin}_{y:\,h(x,y)\le0}g(x,y)}, and \AtlasInlineMath{\Phi(x)=f(x,y^*(x))}. Seek the following approximate stationarity criterion for \AtlasInlineMath{\min_{x\in \mathbb{R}^{d_x}}\Phi(x)}.

\textbf{Status.} Bounds missing

\textbf{TASK tags.} task family: Bilevel

\textbf{MODEL tags.} convexity: Lower objective strongly convex in its lower variable; initialization: Initial distance; initialization: Initial gap

\textbf{ORACLE tags.} interface family: Value gradient

\textbf{ALGORITHMS tags.} restriction: No additional span restriction explicitly imposed

\textbf{OUTPUT tags.} certificate: No additional output certificate required

\textbf{METRIC tags.} criterion: Goldstein stationarity

\textbf{GUARANTEE tags.} aggregation: Per instance

\textbf{RESOURCE tags.} accounting: Joint coordinates retained; measure: N fo f; measure: N fo g

\needspace{7\baselineskip}\section{1. Setting specification}
\needspace{4\baselineskip}\subsection{Task}
Let \AtlasInlineMath{h(x,y)=Ax-By-b}, \AtlasInlineMath{y^*(x)=\operatorname*{argmin}_{y:\,h(x,y)\le0}g(x,y)}, and \AtlasInlineMath{\Phi(x)=f(x,y^*(x))}. Seek the following approximate stationarity criterion for \AtlasInlineMath{\min_{x\in \mathbb{R}^{d_x}}\Phi(x)}.

\needspace{4\baselineskip}\subsection{Admissible model}
The function \AtlasInlineMath{f} has \AtlasInlineMath{C_f}-Lipschitz joint gradient and is \AtlasInlineMath{L_f}-Lipschitz in value. The function \AtlasInlineMath{g} has \AtlasInlineMath{C_g}-Lipschitz joint gradient and is \AtlasInlineMath{\mu_y}-strongly convex in \AtlasInlineMath{y}. All norms are Euclidean or their induced operator norms.

Inputs \AtlasInlineMath{A\in\mathbb{R}^{d_h\times d_x}}, \AtlasInlineMath{B\in\mathbb{R}^{d_h\times d_y}}, and \AtlasInlineMath{b\in\mathbb{R}^{d_h}} are public. The matrix \AtlasInlineMath{B} has full row rank and smallest singular value at least the public \AtlasInlineMath{\beta_B>0}. The initial point has upper-level gap at most \AtlasInlineMath{\Delta} and lower-level initial distance at most \AtlasInlineMath{R_y}.

\AtlasInlineMath{x\in\mathbb{R}^{d_x}}. The solution map \AtlasInlineMath{y^*(\cdot)} is \AtlasInlineMath{L_y}-Lipschitz. The reduced objective is finite and globally Lipschitz, but changes in the active constraints may make it nondifferentiable.

\needspace{4\baselineskip}\subsection{Oracle interface}
Choose function \AtlasInlineMath{f} or \AtlasInlineMath{g} and obtain its exact value and full gradient jointly at an admissible \AtlasInlineMath{(x,y)}. Queries may use the full Euclidean space.

Public linear-constraint operations and corresponding projections are free. Inner solutions, active sets, and hypergradients are not supplied implicitly.

\needspace{4\baselineskip}\subsection{Output}
Return an outer point \AtlasInlineMath{\widehat{x}\in \mathbb{R}^{d_x}}, which may be unqueried. No inner point or certificate is required.

\needspace{4\baselineskip}\subsection{Loss or accuracy criterion}
The error is \AtlasInlineMath{\operatorname{dist}(0,\partial_{r_G}\Phi(\widehat{x}))}, where \AtlasInlineMath{\partial_{r_G}\Phi(x)=\operatorname{conv}\bigcup_{\|z-x\|\le r_G}\partial_C\Phi(z)}; \AtlasInlineMath{\partial_C} denotes the Clarke subdifferential and \AtlasInlineMath{r_G>0} is the public neighborhood radius.

\needspace{4\baselineskip}\subsection{Resource accounting}
Retain the two levels' primitive call counts. A joint value-and-gradient response selects one function and costs one unit; the total primitive cost is \AtlasInlineMath{N=N_{\mathrm{FO},f}+N_{\mathrm{FO},g}}. Every call used in an approximate inner solve is included in the fixed budget.

\needspace{4\baselineskip}\subsection{Algorithm views and success guarantees}
Arbitrary randomized adaptive algorithms are allowed, with measurable queries, stopping, and postprocessing. Outputs may be unqueried, randomly selected historical points, or admissible averages. No span or zero-respecting restriction applies. Queries are sequential; each response in a batch is charged. Internal arithmetic, linear algebra, and storage are free, with no separate work, memory, or communication budget.

For every instance, \AtlasInlineMath{\mathbb{E}\operatorname{dist}(0,\partial_{r_G}\Phi(\widehat{x}))\le\varepsilon}. Expectation is over algorithmic randomness.

\needspace{7\baselineskip}\section{2. Notation and parameters}
\AtlasInlineMath{x,y}: The upper- and lower-level variables. \textbf{Domain:} \AtlasInlineMath{x\in\mathbb R^{d_x},y\in\mathbb R^{d_y}}.

\AtlasInlineMath{f,g}: The original upper- and lower-level real-valued functions. \textbf{Domain:} \AtlasInlineMath{f,g:\mathbb R^{d_x}\times\mathbb R^{d_y}\to\mathbb R}.

\AtlasInlineMath{\widehat x,x_0,y_0}: Output outer point, public outer initial point, and public lower-level initial point \textbf{Domain:} Each lies in its corresponding variable space.

\AtlasInlineMath{\Phi}: The reduced objective. \textbf{Domain:} Defined in the task specification.

\AtlasInlineMath{y^*}: Mapping to the unique lower-level global minimizer \textbf{Domain:} Constraints specified under TASK.

\AtlasInlineMath{\nabla,\|\cdot\|}: Euclidean gradient and vector norm; matrix and higher-derivative norms are the induced Euclidean operator norms

\AtlasInlineMath{\mathbb E,\Pr,\inf,\sup}: Expectation, probability, infimum, and supremum, with aggregation order and randomness specified by the guarantee.

\AtlasInlineMath{\mathcal O,\Omega,\Theta,\widetilde{\mathcal O}}: Asymptotic upper, lower, matching-order, and declared-logarithm upper bounds; fixed parameters are specified in each evidence record

\AtlasInlineMath{c,C}: Positive absolute constants or functions of fixed budgets, with dependencies specified by the corresponding result \textbf{Domain:} \AtlasInlineMath{0<c,C<\infty}.

\AtlasInlineMath{\varepsilon}: Target accuracy. Whether error is squared is determined by the metric and guarantee formulas. \textbf{Domain:} \AtlasInlineMath{\varepsilon>0}.

\AtlasInlineMath{d_x}: Upper-level dimension \textbf{Domain:} Positive integer.

\AtlasInlineMath{d_y}: Lower-level dimension \textbf{Domain:} Positive integer.

\AtlasInlineMath{d_h}: Number of lower-level linear constraints \textbf{Domain:} The quantity \AtlasInlineMath{0\le d_h\le d_y} is an integer; zero denotes no constraints.

\AtlasInlineMath{L_f}: A joint-variable Lipschitz bound on upper-level function values. \textbf{Domain:} \AtlasInlineMath{L_f\ge0}.

\AtlasInlineMath{C_f}: Joint Lipschitz bound for the full upper-level gradient \textbf{Domain:} \AtlasInlineMath{C_f\ge0}.

\AtlasInlineMath{C_g}: Joint Lipschitz bound for the full lower-level gradient \textbf{Domain:} \AtlasInlineMath{C_g\ge\mu_y>0}.

\AtlasInlineMath{\mu_y}: The uniform strong-convexity modulus in the lower-level variable. \textbf{Domain:} \AtlasInlineMath{\mu_y>0}.

\AtlasInlineMath{\Delta}: Public upper bound on the initial reduced-objective gap \textbf{Domain:} \AtlasInlineMath{\Delta\ge0}.

\AtlasInlineMath{\Phi(x_0)-\inf_x\Phi(x)\le\Delta}

\AtlasInlineMath{R_y}: Public initial lower-level distance bound \textbf{Domain:} \AtlasInlineMath{R_y\ge0}.

\AtlasInlineMath{\|y_0-y^*(x_0)\|\le R_y}

\AtlasInlineMath{\beta_B}: A lower bound on the smallest singular value of the public constraint matrix B. \textbf{Domain:} \AtlasInlineMath{\beta_B>0}.

\AtlasInlineMath{L_y}: A global Lipschitz bound on the lower-level solution map. \textbf{Domain:} \AtlasInlineMath{L_y\ge0}.

\AtlasInlineMath{\|y^*(x)-y^*(x\prime)\|\le L_y\|x-x\prime\|}

\AtlasInlineMath{r_G}: The neighborhood radius in the Goldstein subdifferential definition. \textbf{Domain:} \AtlasInlineMath{r_G>0}.

\AtlasInlineMath{\kappa_g}: Lower-level condition number \textbf{Domain:} \AtlasInlineMath{\kappa_g\ge1}.

\AtlasInlineMath{\kappa_g=C_g/\mu_y}

\AtlasInlineMath{A,B,b}: The public lower-level linear-constraint coefficients. \textbf{Domain:} \AtlasInlineMath{A\in\mathbb R^{d_h\times d_x},B\in\mathbb R^{d_h\times d_y},b\in\mathbb R^{d_h}}.

\AtlasInlineMath{h}: Residual of the public linear constraints

\AtlasInlineMath{h(x,y)=Ax-By-b}

\AtlasInlineMath{\partial_C,\partial_{r_G},\operatorname{conv},\operatorname{dist}}: Clarke subdifferential, Goldstein subdifferential, convex hull, and Euclidean point-to-set distance; full definitions appear under METRIC

\AtlasInlineMath{L_0}: A Lipschitz bound for reduced-objective values derived from the original-function and solution-map budgets.

\AtlasInlineMath{L_0=L_f\sqrt{1+L_y^2}}

\AtlasInlineMath{\alpha}: Algorithm-selected accuracy for approximating the reduced objective value \textbf{Domain:} Require \AtlasInlineMath{0<\alpha\le\varepsilon/2} and the internal-parameter conditions of the stated upper bound.

\AtlasInlineMath{N,Q}: This is the total call budget under the stated resource projection. If only one symbol is used, the two names are aliases for that same budget. \textbf{Domain:} RESOURCE specifies which operations are summed; no additional unit cost is introduced.

\AtlasInlineMath{N_{\rm FO,f},N_{\rm FO,g}}: Counts of joint value and full-gradient responses for separately selected upper- or lower-level functions. \textbf{Domain:} Nonnegative integer.

\needspace{12\baselineskip}\section{3. Complexity bracket}
\needspace{4\baselineskip}\subsection{Compared quantities}
\needspace{4\baselineskip}\subsubsection{Q1 - target accuracy cost}
\AtlasMath{Q^*(\varepsilon)=\inf_A\sup_{I,\mathcal O}N(A,I,\mathcal O)}
\needspace{4\baselineskip}\subsection{R1 - The high-accuracy Goldstein region.}
\textbf{Region.} Assume \AtlasInlineMath{r_G>0} and \AtlasInlineMath{\Delta>0}, with \AtlasInlineMath{\varepsilon} sufficiently small. Set inner function-value accuracy to \AtlasInlineMath{0<\alpha\le\varepsilon/2} and the source outer tolerance to \AtlasInlineMath{\varepsilon/2}. All other model budgets are fixed.

\textbf{Quantity.} Q1

\textbf{Resource.} Joint resource budget

\textbf{Status.} Lower bound missing

\textbf{Lower bound.} No admitted bound for this comparison

\textbf{Upper bound.} B1

\textbf{Gap or comparison.} Upper bounds using the original interface are included. Results requiring exact multipliers or active sets are not mixed into the profile.

\needspace{7\baselineskip}\section{4. Bounds and supporting evidence}
The following conditions apply to every bound in this profile, in addition to each bound's stated region and algorithm scope.

\textbf{Conditions.} The algorithm selects the internal tolerance \AtlasInlineMath{\alpha} so that \AtlasInlineMath{L_0^2} dominates the bracket. This gives \AtlasInlineMath{\widetilde O(\sqrt{\kappa_g}d_x\Delta L_0^2/(r_G\varepsilon^3))} for nondegenerate budgets.

\textbf{Dimension conditions.} The upper bound holds for every finite-dimensional instance in the declared model. Any dimension or fixed-budget dependence not displayed explicitly follows the accompanying evidence qualification.

\needspace{4\baselineskip}\subsection{B1 - Approximate reduced values, random-direction smoothing, and an online-to-nonconvex conversion.}
\textbf{Admission.} Included in the bracket

\textbf{Direction.} upper

\textbf{Witness or lower-bound scope.} Randomized witness

\AtlasMath{\mathcal{O}\!\left(\frac{\sqrt{\kappa_g}\,d_x\Delta}{r_G\varepsilon^3}\left[L_0^2+\alpha^2\left(\frac{d_x}{r_G^2}+\frac{d_x^2L_0^2}{\Delta^2}\right)\right]\log\frac{L_f}{\alpha}\right)}
\textbf{Sources.} \href{https://proceedings.neurips.cc/paper_files/paper/2024/file/ffd86e56e6403d63dd6face033060e5a-Paper-Conference.pdf}{[1]}

\textbf{Constants.} Unshown regularity, initialization, and noise budgets are fixed. Their joint parameter dependence is not certified as closed.

\textbf{Derivation.} Use the conservative reduced-objective Lipschitz bound \AtlasInlineMath{L_0=L_f\sqrt{1+L_y^2}} in the general Lipschitz analysis of Theorem C.1. Retaining the direct variation of \AtlasInlineMath{f} in \AtlasInlineMath{x} covers \AtlasInlineMath{L_y=0}.

\textbf{Source location.} \href{https://proceedings.neurips.cc/paper_files/paper/2024/file/ffd86e56e6403d63dd6face033060e5a-Paper-Conference.pdf}{[1]} theorem: Theorem 4.1; section: 4

The source total error is \AtlasInlineMath{\varepsilon_{\rm out}+\alpha}. Set \AtlasInlineMath{\varepsilon_{\rm out}=\varepsilon/2} and \AtlasInlineMath{\alpha\le\varepsilon/2}. Absorb the factor 2 into the big-O constant while retaining the source expression in which \AtlasInlineMath{\alpha^2} multiplies the entire bracket.

\needspace{7\baselineskip}\section{References}
[1] Guy Kornowski; Swati Padmanabhan; Kai Wang; Jimmy Zhang; Suvrit Sra. First-Order Methods for Linearly Constrained Bilevel Optimization. 2024. NeurIPS2024 proceedings PDF.

\url{https://proceedings.neurips.cc/paper_files/paper/2024/file/ffd86e56e6403d63dd6face033060e5a-Paper-Conference.pdf}

\end{document}
